\documentclass{article}

\usepackage[english]{babel}
\usepackage[letterpaper,top=2cm,bottom=2cm,left=3cm,right=3cm,marginparwidth=1.75cm]{geometry}

\usepackage{amsmath}
\usepackage{graphicx}
\usepackage{xcolor}
\usepackage[colorlinks=true, allcolors=blue]{hyperref}
\usepackage{subcaption}
\usepackage{tabularray}
\usepackage{float}
\usepackage[normalem]{ulem}

\UseTblrLibrary{booktabs}
\UseTblrLibrary{siunitx}

\newcommand{\tinytableTabularrayUnderline}[1]{\underline{#1}}
\newcommand{\tinytableTabularrayStrikeout}[1]{\sout{#1}}
\NewTableCommand{\tinytableDefineColor}[3]{\definecolor{#1}{#2}{#3}}

\title{Problem Set 10}
\date{April 21, 2026}
\author{Caleb Reid}

\begin{document}

\maketitle

For the logit model, the penalty, or lambda, was 0.0000000001. For all intensive purposes, the penalty parameter is essentially zero.

For the tree model, the ideal cost complexity is 0.001, which is the minimum penalty value possible. The ideal tree depth is 15 and the minimum sample size for making a split is 10, which again is the minimum penalty value possible.

For the K-Nearest Neighbors model, the ideal number of neighbors is 30, which is the maximum number of neighbors possible.

For the neural networks model, the ideal number of hidden units is seven and the penalty parameter is one.

In terms of accuracy, the tree model had the highest accuracy at 86.8\% (estimate = 0.868417). In second place, the logit model reported an accuracy estimate of 0.852602 or 85.3\%. The neural networks and k-nearest neighbors model reported similar accuracies of 0.843697 (84.4\%) and 0.84339 (84.3\%), respectively.  

\end{document}